\chapter{DISTRIBUSI}

\section{Limit Induktif}\label{sec_ind_limits}
\begin{defn}\label{defn_dir_sys} Misalkan $D$ suatu himpunan terarah terhadap pengurutan parsial~$\le$ dan misalkan $\{A_i\colon i \in D\}$ suatu keluarga
objek dalam suatu kategori~$\cat C$.  Anggap bahwa untuk setiap $i$, $j \in D$ dengan $i \le j$ terdapat suatu morfisme
$\phi_{j,i}\colon A_i \sto A_j$ yang memenuhi $\phi_{ki} = \phi_{kj}\phi_{ji}$ setiap kali $i \le j \le k$ dalam~$D$.  Anggap pula bahwa
$\phi_{ii}$ adalah pemetaan identitas pada $A_i$ untuk setiap $i \in D$. Maka keluarga berindeks $\vc A = \bigl(A_i\bigr)_{i \in D}$ dari
objek-objek bersama keluarga berindeks $\bs\phi = \bigl(\phi_{ji}\bigr)$ dari \emph{morfisme penghubung} disebut suatu
 \index{morfisme penghubung}%
 \index{morfisme!penghubung}%
 \index{terarah!sistem}%
 \index{sistem!terarah}%
\df{sistem terarah} dalam~$\cat C$.
\end{defn}

\begin{defn}\label{defn_ind_lim} Misalkan pasangan $(\vc A, \bs\phi)$ suatu sistem terarah (seperti di atas) dalam kategori~$\cat C$ yang
himpunan terarah dasarnya adalah~$D$.  Suatu
 \index{induktif!limit}%
 \index{limit!induktif}%
\df{limit induktif} (atau
 \index{langsung!limit}%
 \index{limit!langsung}%
\df{limit langsung}) dari sistem $(\vc A,\bs\phi)$ adalah pasangan $(L,\bs\mu)$ dengan $L$ suatu objek dalam $\cat C$
dan $\bs\mu = \bigl(\mu_i\bigr)_{i \in D}$ suatu keluarga morfisme $\mu_j\colon A_j \sto L$ dalam $\cat C$ yang memenuhi
 \begin{enumerate}[label=\rm{(\alph*)}]
  \item $\mu_i = \mu_j \phi_{ji}$ setiap kali $i \le j$ dalam $D$, dan
  \item jika $(M,\bs\lambda)$ suatu pasangan dengan $M$ suatu objek dalam $\cat C$ dan $\bs\lambda = \bigl(\lambda_i\bigr)_{i \in D}$
adalah suatu keluarga berindeks morfisme $\lambda_j\colon A_j \sto M$ dalam $\cat C$ yang memenuhi $\lambda_i = \lambda_j\phi_{ji}$ setiap kali $i \le j$ dalam~$D$,
maka terdapat morfisme tunggal $\psi\colon L \sto M$ sedemikian sehingga $\lambda_i = \psi\mu_i$ untuk setiap $i \in D$.
 \end{enumerate}
Dengan penyalahgunaan bahasa yang lazim, biasanya kita mengatakan bahwa $L$ adalah limit induktif dari sistem $\vc A = (A_i)$ dan
 \index{<arrowto@$\underrightarrow{\lim} A_i$ (limit induktif)}%
menulis $L = \underrightarrow{\lim} A_i$.

\vskip 15 pt

   \[\xymatrix@+15pt@C+25pt{&&&& L\ar@{-->}[dddd]^\psi \\
                   &&&&   \\
                  {\cdots}\ar[r] & A_i \ar[uurrr]^{\mu_i}\ar[ddrrr]^{\lambda_i}\ar[r]^(.6){\phi_{ji}} &
                       A_j\ar[uurr]^{\mu_j}\ar[ddrr]^{\lambda_j}\ar[r] & {\cdots} &  \\
                   &&&&   \\
                   &&&& M
    }\]
\end{defn}

Dalam suatu kategori (konkret) sebarang, limit induktif belum tentu ada.  Namun, jika ada, limit tersebut tunggal.

\begin{prop}\label{prop_ind_lim_unique} Limit induktif (jika ada dalam suatu kategori) bersifat tunggal (hingga isomorfisme).
\end{prop}

\begin{exam} Misalkan $S$ suatu objek tak kosong dalam kategori $\cat{SET}$. Tinjau keluarga $\sfml P(S)$ dari semua himpunan bagian $S$ yang diarahkan
oleh pengurutan parsial~$\subseteq$.  Setiap kali $A \subseteq B \subseteq S$, ambil morfisme penghubung $\sbsb\iota{BA}$ sebagai
pemetaan inklusi dari $A$ ke~$B$. Maka $\sfml P(S)$, bersama keluarga morfisme penghubung, membentuk suatu sistem terarah.  Limit
induktif dari sistem ini adalah~$S$.
\end{exam}

\begin{exam} Limit langsung ada dalam kategori $\cat{VEC}$ dari ruang vektor dan pemetaan linear.
\end{exam}

\begin{proof}[\emph{Petunjuk untuk bukti}] Misalkan $(\vc V, \bs\phi)$ suatu sistem terarah dalam $\cat{VEC}$ (seperti dalam definisi~\ref{defn_dir_sys})
dan $D$ adalah himpunan terarah dasarnya. Misalkan
$W=\bigoplus_{i\in D}V_i$ dan $\iota_i\colon V_i\sto W$ inklusi kanonik. Tetapkan
  \[ N=\operatorname{span}\{\iota_i(u)-\iota_j(\phi_{ji}(u)):
       i\le j,\quad u\in V_i\}. \]
Ambil $L=W/N$ dan $\mu_i=q\circ\iota_i$, dengan $q\colon W\sto W/N$ pemetaan hasil bagi, lalu verifikasikan sifat universalnya
(lihat contoh~\ref{C015544} dan~\ref{X_dir_sum_VEC}).  \ns
\end{proof}

\begin{exam} Barisan $\Z \to^{\vc 1} \Z \to^{\vc 2} \Z \to^{\vc 3} \Z \to^{\vc 4} \dots$
(dengan $\vc n \colon \Z \sto \Z$ memenuhi $\vc n(1) = n$) merupakan suatu barisan induktif grup-grup Abelian
yang limit induktifnya adalah himpunan $\Q$ dari bilangan rasional.
\end{exam}

\begin{exam} Barisan $\Z \to^{\vc 2} \Z \to^{\vc 2} \Z \to^{\vc 2} \Z \to^{\vc 2} \dots$
merupakan suatu barisan induktif grup-grup Abelian yang limit induktifnya adalah himpunan bilangan rasional
diadik.
\end{exam}









\section{Ruang-$LF$}\label{sec_LFsps}
Karena tujuan utama kita dalam bab ini adalah pengantar teori distribusi, mulai sekarang kita akan membatasi perhatian
pada jenis limit induktif yang sangat sederhana---\emph{limit induktif ketat dari suatu barisan induktif}. Secara khusus, kita akan
tertarik pada limit induktif ketat dari barisan ruang Fr\'echet yang meningkat.


Ingat bahwa topologi lemah pada suatu himpunan $S$ diinduksi oleh suatu keluarga fungsi $f_\alpha \colon S \sto
X_\alpha$ yang memetakan ke ruang-ruang topologis. Topologi ini adalah topologi terlemah pada $S$ yang membuat
semua fungsi tersebut kontinu (lihat latihan~\ref{C021421}). Gagasan dualnya adalah
 \index{topologi!kuat}%
 \index{kuat!topologi}%
\df{topologi kuat}, yakni topologi yang diinduksi oleh suatu keluarga fungsi yang memetakan \emph{dari}
ruang-ruang topologis $X_\alpha$ \emph{ke dalam}~$S$. Topologi ini adalah topologi terkuat yang membuat
semua fungsi tersebut kontinu. Salah satu topologi kuat yang telah kita jumpai adalah
\emph{topologi hasil bagi} (lihat proposisi~\ref{prop_quotient_top_strong}). Contoh penting lainnya,
yang akan kita perkenalkan berikutnya, adalah \emph{topologi limit induktif}.


\begin{defn} Suatu
 \index{induktif!barisan!ketat}%
 \index{barisan!induktif ketat}%
 \index{ketat!barisan induktif}%
\df{barisan induktif ketat} adalah barisan $(X_i)$ ruang konveks lokal sedemikian sehingga untuk setiap $i \in \N$
   \begin{enumerate}
     \item[(i)] $X_i$ adalah subruang tertutup dari $X_{i+1}$ dan
     \item[(ii)] topologi $\sfml T_i$ pada $X_i$ adalah pembatasan topologi $\sfml T_{i+1}$ pada $X_i$.
   \end{enumerate}
Ini, tentu saja, merupakan sistem terarah dengan morfisme penghubung $\phi_{ji}$ yang tidak lain adalah pemetaan inklusi dari $X_i$ ke $X_j$ untuk $i < j$.

Yang disebut
 \index{induktif!limit!ketat}%
 \index{limit!induktif ketat}%
 \index{ketat!limit induktif}%
\df{limit induktif ketat} dari barisan ruang semacam itu adalah gabungannya $L = \bigcup_{i=1}^\infty X_i$. Kita memberi $L$ topologi
konveks lokal terbesar yang membuat semua pemetaan inklusi $\psi_i\colon X_i \sto L$ kontinu. Inilah
 \index{induktif!limit!topologi}%
 \index{topologi!limit induktif}%
 \index{limit!induktif!topologi}%
\df{topologi limit induktif} pada~$L$.
\end{defn}

\begin{prop} Topologi limit induktif yang didefinisikan di atas ada.
\end{prop}

\begin{defn} Limit induktif ketat dari suatu barisan ruang Fr\'echet disebut
 \index{LF@ruang-$LF$}%
 \index{ruang!$LF$-}%
\df{ruang-$LF$}.
\end{defn}

\begin{exam} Ruang $\fml D(\Omega) = \fml C^\infty_c(\Omega)$ dari fungsi uji pada suatu himpunan terbuka $\Omega$ dalam ruang Euklides
adalah ruang-$LF$ (lihat notasi~\ref{mi_notn}).
\end{exam}

\begin{prop} Setiap ruang-$LF$ lengkap.
\end{prop}

\begin{proof} Lihat~\cite{Treves:1967}, Teorema 13.1.  \ns
\end{proof}

\begin{prop} Misalkan $T\colon L \sto Y$ suatu pemetaan linear dari ruang-$LF$ ke ruang konveks lokal. Andaikan bahwa $L = \underrightarrow{\lim}X_k$
dengan $(X_k)$ suatu barisan induktif ketat ruang Fr\'echet. Maka $T$ kontinu jika dan hanya jika pembatasannya $T\bigr|_{X_k}$
pada setiap $X_k$ kontinu.
\end{prop}

\begin{proof} Lihat~\cite{Treves:1967}, Proposisi~13.1, atau~\cite{Conway:1990}, Proposisi IV.5.7, atau~\cite{Grubb:2009}, Teorema B.18(c).  \ns
\end{proof}

\begin{cau} Fran\c{c}ois Treves, dalam bukunya yang ditulis dengan indah, \emph{Topological Vector Spaces, Distributions, and Kernels}, memperingatkan
para pembacanya tentang suatu kekeliruan yang sangat sulit dikenali: orang mudah sekali menganggap bahwa jika $M$ adalah subruang vektor tertutup dari ruang-$LF$
$L = \underrightarrow{\lim}X_k$, maka topologi yang diwarisi $M$ dari $L$ pasti sama dengan topologi limit induktif yang diperolehnya
dari barisan induktif ketat ruang Fr\'echet~$M \cap X_k$. Ternyata hal ini \emph{tidak} benar secara umum.
\end{cau}

\begin{prop} Misalkan $T\colon L \sto Y$ suatu pemetaan linear dari ruang-$LF$ ke ruang konveks lokal. Maka $T$ kontinu jika dan hanya jika
pemetaan tersebut kontinu sekuensial.
\end{prop}

















\section{Distribusi}
\begin{defn} Misalkan $\Omega$ suatu himpunan bagian terbuka tak kosong dari~$\R^n$. Suatu fungsi $f\colon \Omega \sto \R$ disebut
 \index{terintegralkan!secara lokal}%
 \index{secara lokal!terintegralkan}%
\df{terintegralkan secara lokal} jika $\int_K \abs{f}\,d\mu < \infty$ untuk setiap himpunan bagian kompak $K$ dari~$\Omega$. (Di sini $\mu$ adalah ukuran Lebesgue
berdimensi-$n$ yang biasa.) Keluarga semua fungsi tersebut pada $\Omega$ dinotasikan
 \index{L@$\locint\Omega$!fungsi yang terintegralkan secara lokal pada~$\Omega$}%
dengan~$\locint\Omega$.
\end{defn}

\begin{exam} Fungsi konstan $\vc 1$ pada garis real terintegralkan secara lokal, meskipun fungsi ini jelas tidak terintegralkan.
\end{exam}

Bahkan fungsi yang terintegralkan secara lokal tidak harus kontinu. Pada garis real, misalnya, fungsi karakteristik dari sebarang himpunan terukur Lebesgue
terintegralkan secara lokal. Bayangkan betapa memudahkannya jika kita dapat membicarakan secara bermakna \emph{turunan} dari sebarang fungsi demikian dan,
lebih baik lagi, diperbolehkan menurunkannya sesering yang kita inginkan. --- Selamat datang di dunia distribusi. Di sini
   \begin{enumerate}
       \item setiap fungsi yang terintegralkan secara lokal dapat dipandang sebagai suatu distribusi, dan
       \item setiap distribusi dapat diturunkan tak berhingga kali, dan lebih lanjut,
       \item semua aturan diferensiasi yang lazim dalam kalkulus dasar berlaku.
   \end{enumerate}

\begin{defn} Misalkan $\Omega$ suatu himpunan bagian terbuka tak kosong dari~$\R^n$. Suatu
 \index{distribusi}%
\df{distribusi} adalah fungsional linear kontinu pada ruang fungsi uji $\fml D(\Omega) = \fml C^\infty_c(\Omega)$ pada~$\Omega$.
Meskipun secara teknis fungsional demikian didefinisikan pada suatu ruang fungsi pada $\Omega$, kita akan mengikuti kebiasaan longgar yang lazim
dan menyebutnya sebagai \emph{distribusi pada~$\Omega$}. Himpunan semua distribusi pada $\Omega$, yaitu ruang dual dari $\fml D(\Omega)$,
 \index{D@$\fml D^*(\Omega)$ (ruang distribusi pada~$\Omega$)}%
akan dinotasikan dengan~$\fml D^*(\Omega)$.
\end{defn}

\begin{prop} Misalkan $\Omega$ suatu himpunan bagian terbuka dari~$\R^n$. Suatu fungsional linear $L$ pada ruang fungsi uji $\fml D(\Omega)$ pada $\Omega$ adalah
distribusi jika dan hanya jika $L(\phi_k) \sto 0$ setiap kali $(\phi_k)$ merupakan barisan fungsi uji sedemikian sehingga $\phi_k^{(\bs\alpha)} \sto 0$
secara seragam untuk setiap multi-indeks $\bs\alpha$ dan terdapat suatu himpunan kompak $K \subseteq \Omega$ yang memuat tumpuan setiap~$\phi_k$.
\end{prop}

\begin{proof} Lihat~\cite{Conway:1990}, proposisi IV.5.21 atau~\cite{Treves:1967}, proposisi 21.1(b).  \ns
\end{proof}

\begin{exam} Jika $\Omega$ suatu himpunan bagian terbuka tak kosong dari $\R^n$ dan $f$ suatu fungsi yang terintegralkan secara lokal pada $\Omega$, maka pemetaan
   \[ L_f\colon \fml D(\Omega) \sto \K \colon \phi \mapsto \int f\phi\,d\lambda \]
merupakan suatu distribusi. Distribusi demikian disebut
 \index{regular!distribusi}%
 \index{distribusi!regular}%
 \index{L@$L_f$ (distribusi regular yang bersesuaian dengan~$f$)}%
distribusi \df{regular}. Distribusi yang tidak regular disebut
 \index{singular!distribusi}%
 \index{distribusi!singular}%
distribusi \df{singular}.
\end{exam}

\emph{Catatan.} Misalkan $u$ suatu distribusi dan $\phi$ suatu fungsi uji pada $\Omega$. Notasi baku untuk $u(\phi)$ adalah
$\langle u,\phi \rangle$ dan $\langle \phi,u \rangle$. Kita juga menggunakan
 \index{<unoptop@$\tilde f$ (distribusi regular yang bersesuaian dengan~$f$)}%
notasi $\tilde f$ untuk $L_f$. Jadi, dalam contoh sebelumnya,

    \[ \langle \tilde f,\phi \rangle = \langle \phi,\tilde f \rangle
            = \tilde f(\phi) = \langle L_f,\phi \rangle
            = \langle \phi,L_f \rangle = L_f(\phi) \]
untuk $f \in \locint\Omega$.

\begin{exam} Jika $\mu$ suatu ukuran Borel regular pada himpunan bagian terbuka tak kosong $\Omega$ dari~$\R^n$ dan $f$ suatu fungsi yang terintegralkan secara lokal
terhadap ukuran tersebut pada $\Omega$, maka pemetaan
    \[ L_\mu\colon \fml D(\Omega) \sto \R\colon \phi \mapsto \int_\Omega f\phi \,d\mu \]
merupakan suatu distribusi pada $\Omega$.
\end{exam}

\begin{notn} Misalkan $a \in \R$. Definisikan $\mu_a$, yaitu
 \index{Dirac!ukuran}%
 \index{ukuran!Dirac}%
 \index{mua@$\mu_a$ (ukuran Dirac yang terkonsentrasi di~$a$)}%
\df{ukuran Dirac} yang terkonsentrasi di $a$, pada himpunan-himpunan Borel dari $\R$, dengan
   \[ \mu_a(B) = \begin{cases}  1, &\text{jika $a \in B$} \\
                                0, &\text{jika $a \notin B$}.
                 \end{cases} \]
Distribusi yang bersesuaian kita notasikan dengan $\delta_a$ dan kita sebut
 \index{Dirac!distribusi}%
 \index{distribusi!Dirac}%
 \index{delta@$\delta_a$, $\delta$ (distribusi Dirac)}%
\df{distribusi delta Dirac di $a$}. Jadi, $\delta_a = L_{\mu_a}$. Jika $a = 0$, kita menulis $\delta$ untuk $\delta_0$.
Secara historis, distribusi $\delta$ telah menyandang nama yang terkenal menyesatkan (meskipun secara teknis
tidak salah): \emph{fungsi delta Dirac}.
\end{notn}

\begin{exam} Jika $\delta_a$ adalah \emph{distribusi delta Dirac} di suatu titik $a \in \R$ dan $\phi$ suatu fungsi uji pada~$\R$,
maka $\delta_a(\phi) = \phi(a)$.
\end{exam}

\begin{notn} Misalkan $H$ menyatakan fungsi karakteristik interval $[0,\infty)$. Fungsi ini adalah
 \index{Heaviside!fungsi}%
 \index{fungsi!Heaviside}%
\df{fungsi Heaviside}. Distribusi yang bersesuaian, $\widetilde H = L_H$, adalah
 \index{Heaviside!distribusi}%
 \index{distribusi!Heaviside}%
 \index{H@$\wt H$ (distribusi Heaviside)}%
\df{distribusi Heaviside}.
\end{notn}

\begin{exam} Jika $\wt H$ adalah \emph{distribusi Heaviside} dan $\phi \in \fml D(\R)$, maka $\wt H(\phi) = \int_0^\infty \phi$.
\end{exam}

\begin{exam} Misalkan $f$ suatu fungsi yang dapat didiferensialkan pada $\R$ dan turunannya terintegralkan secara lokal. Maka
    \begin{equation}\label{deriv_reg_distr}
        L_{f\,'}(\phi) = -L_f(\phi\,')
    \end{equation}
untuk setiap $\phi \in \fml D(\R)$.
\end{exam}

\begin{proof}[\emph{Petunjuk untuk bukti}] Integrasi parsial. \ns
\end{proof}

Sekarang timbul sebuah pertanyaan penting: Bagaimana kita dapat mendefinisikan ``turunan'' suatu distribusi? Tentunya wajar jika kita menginginkan suatu
distribusi regular berperilaku hampir sama dengan fungsi yang melahirkannya. Artinya, kita ingin turunan $(L_f)'$ dari
suatu distribusi regular $L_f$ bersesuaian dengan distribusi regular yang timbul dari turunan fungsi~$f$. Dengan kata lain,
yang kita inginkan adalah $(L_f)' = L_{f\,'}$. Menurut contoh sebelumnya, ini menyarankan bahwa untuk fungsi $f$ yang dapat didiferensialkan
yang turunannya terintegralkan secara lokal, kita seharusnya mendefinisikan
   \begin{equation}\label{deriv_distr1}
        (L_f)'(\phi) = -L_f(\phi\,')
   \end{equation}
untuk setiap $\phi \in \fml D(\R)$. Pengamatan bahwa ruas kanan persamaan~\eqref{deriv_distr1} bermakna untuk
\emph{distribusi apa pun} memotivasi definisi berikut.

\begin{defn}\label{def_deriv_distr} Misalkan $u$ suatu distribusi pada himpunan bagian terbuka tak kosong $\Omega$ dari~$\R$. Kita mendefinisikan $u\,'$, yaitu
 \index{turunan!distribusi}%
 \index{distribusi!turunan}%
\df{turunan} dari $u$,
 \index{D@$Du$ (turunan suatu distribusi~$u$)}%
 \index{<unopur@$u'$ (turunan suatu distribusi~$u$)}%
dengan
    \[ u\,'(\phi) := -u(\phi\,') \]
untuk setiap fungsi uji $\phi$ pada~$\Omega$. Kita juga menggunakan notasi $Du$ untuk turunan dari~$u$.
\end{defn}

\begin{exam} \emph{Distribusi Dirac} $\delta$ adalah turunan dari \emph{distribusi Heaviside}~$\widetilde H$.
\end{exam}

\begin{exer} Misalkan $S$ suatu fungsi signum (atau tanda) pada $\R$; artinya, misalkan $S = 2H - 1$, dengan $H$ fungsi Heaviside.
Misalkan pula $g$ suatu fungsi pada $\R$ sedemikian sehingga $g\colon x \mapsto -x - 3$ untuk $x < 0$ dan $g\colon x \mapsto x + 5$ untuk $x \ge 0$.
   \begin{enumerate}[label=\rm{(\alph*)}]
     \item Carilah konstanta $\alpha$ dan $\beta$ sedemikian sehingga ${\tilde g}' = \alpha \widetilde S + \beta \delta$.
     \item Carilah suatu fungsi $a$ pada $\R$ sedemikian sehingga ${\tilde a}' = \wt S$.
   \end{enumerate}
\end{exer}

\begin{exer} Misalkan $f(x) = \left\{
                           \begin{array}{ll}
                             1 - \abs x, & \hbox{jika $\abs x \le 1$;} \\
                                 0,      & \hbox{selain itu.}
                           \end{array}
                         \right.$ Dengan menghitung kedua ruas persamaan~\eqref{deriv_reg_distr} secara terpisah, tunjukkan bahwa persamaan itu
benar secara klasik untuk setiap fungsi uji yang tumpuannya termuat dalam interval~$[-1,1]$.
\end{exer}

\begin{exer} Misalkan $f(x) = \sbsb\chi{(0,1)}$ untuk $x \in~\R$ dan $\phi \in \fml D\bigl((-1,1)\bigr)$. Bandingkan nilai $L_{f\,'}(\phi)$ dan
$\bigl(L_f\bigr)'(\phi)$.
\end{exer}

\begin{exer} Misalkan $f(x) = \left\{
                           \begin{array}{ll}
                             x, & \hbox{jika $0 \le x \le 1$;} \\
                             x + 2, & \hbox{jika $1 < x \le 2$;} \\
                             0, & \hbox{selain itu.}
                           \end{array}
                         \right.$ Bandingkan nilai $L_{f\,'}(\phi)$ dan $\bigl(L_f\bigr)'(\phi)$ untuk suatu fungsi uji~$\phi$.
\end{exer}

\begin{exer} Misalkan $h(x) = \left\{
                           \begin{array}{ll}
                             2x, & \hbox{jika $x < -1$;} \\
                             1, & \hbox{jika $-1 \le x < 1$;} \\
                             0, & \hbox{jika $x \ge 1$.}
                           \end{array}
                         \right.$ Carilah suatu fungsi $f$ yang terintegralkan secara lokal sedemikian sehingga ${\tilde f}\,' = \tilde h + 3\delta_1 - \delta_{-1}$.
\end{exer}

Definisi~\ref{def_deriv_distr} dapat diperluas ke turunan parsial dari distribusi tanpa kesulitan.

\begin{defn} Misalkan $u$ suatu distribusi pada suatu himpunan bagian terbuka tak kosong $\Omega$ dari $\R^n$ untuk suatu $n \ge 2$ dan $\bs\alpha$
suatu multi-indeks. Kita mendefinisikan
 \index{operator diferensial!yang bekerja pada distribusi}%
 \index{D@$D^{\bs\alpha}$ (notasi multi-indeks untuk diferensiasi)}%
\df{operator diferensial} $D^{\bs\alpha} = \bigl(\frac{\partial}{\partial x_1}\bigr)^{\alpha_1} \dots
\bigl(\frac{\partial}{\partial x_n}\bigr)^{\alpha_n}$ berorde $\abs{\bs\alpha}$ yang bekerja pada $u$ dengan
   \[ D^{\bs\alpha}u(\phi) := (-1)^{\abs{\bs\alpha}}u\bigl(\phi^{(\bs\alpha)}) \]
untuk setiap fungsi uji $\phi$ pada~$\Omega$.
 \index{<unopura@$u^{\bs\alpha}$ (notasi multi-indeks untuk diferensiasi)}%
Notasi alternatif untuk $D^{\bs\alpha}u$ adalah $u^{(\bs\alpha)}$.
\end{defn}

\begin{exer} Suatu \df{dipol} (dengan momen listrik 1 di titik asal pada $\R$) dapat dipandang sebagai ``limit'' ketika $\epsilon \sto 0^+$
 \index{dipol}%
dari suatu sistem $T_\epsilon$ yang terdiri atas dua muatan $-\frac1\epsilon$ dan $\frac1\epsilon$ yang masing-masing ditempatkan di $0$ dan $\epsilon$. Tunjukkan
bagaimana hal ini dapat mengarahkan kita untuk mendefinisikan dipol secara matematis sebagai $-\delta\,'$. \emph{Petunjuk.} Pandanglah $T_\epsilon$ sebagai
distribusi $\frac1\epsilon \delta_\epsilon - \frac1\epsilon \delta$.
\end{exer}
 
\begin{defn} Misalkan $X$ suatu ruang konveks lokal Hausdorff dan $X^*$ ruang dualnya, yaitu ruang semua
fungsional linear kontinu pada~$X$. Untuk setiap $f \in X^*$ definisikan
  \[ p_f\colon X \sto \K\colon x \mapsto \abs{f(x)}. \]
Jelas bahwa setiap $p_f$ merupakan seminorma pada~$X$. Seperti dalam kasus ruang linear bernorma, kita mendefinisikan
 \index{lemah!topologi!pada ruang konveks lokal}%
 \index{sigma@$\sigma(X,X^*)$ (topologi lemah pada~$X$)}%
\df{topologi lemah} pada $X$, yang dinotasikan dengan $\sigma(X,X^*)$, sebagai topologi lemah yang dibangkitkan oleh keluarga seminorma
$\{p_f\colon f \in X^*\}$ pada~$X$. Demikian pula, untuk setiap $x \in X$ definisikan
   \[ p_x\colon X^* \sto \K \colon f \mapsto \abs{f(x)}. \]
Sekali lagi, setiap $p_x$ merupakan seminorma pada~$X^*$, dan kita mendefinisikan
 \index{w*@$w^*$-topologi (topologi lemah-bintang)}%
 \index{sigma@$\sigma(X^*,X)$ (topologi-$w^*$ pada~$X^*$)}%
\df{topologi-$w^*$} pada $X^*$, yang dinotasikan dengan $\sigma(X^*,X)$, sebagai topologi lemah yang dibangkitkan oleh keluarga seminorma
$\{p_x\colon x \in X\}$ pada~$X^*$.
\end{defn}

\begin{prop} Jika ruang distribusi $\fml D^*(\Omega)$ pada subhimpunan terbuka $\Omega$ dari $\R^n$ diberi topologi-$w^*$,
maka suatu jaring $(u_\lambda)$ distribusi pada $\Omega$ konvergen ke distribusi $v$ pada $\Omega$ jika dan hanya jika
$\left\langle u_\lambda, \phi\right\rangle \sto \left\langle v, \phi\right\rangle$ untuk setiap fungsi uji $\phi$ pada $\Omega$.
\end{prop}

\begin{prop} Andaikan suatu jaring $(u_\lambda)$ distribusi konvergen ke distribusi~$v$. Maka ${u_\lambda}^{(\alpha)} \sto v^{(\alpha)}$
untuk setiap multi-indeks $\alpha$.
\end{prop}

\begin{proof} Lihat \cite{Rudin:1991}, Teorema 6.17; atau \cite{Yosida:1965}, Bab II, Seksi 3, Teorema; atau \cite{Grubb:2009}, Teorema 3.9;
atau~\cite{HirschL:1999}, Bab 8, Proposisi 2.4.
\end{proof}

\begin{prop} Jika $g$, $f_k \in \locint\Omega$ untuk setiap $k$ (dengan $\open{\Omega}{\R^n}$) dan $f_k \sto g$ secara seragam pada
setiap subhimpunan kompak dari $\Omega$, maka $\tilde f_k \sto \tilde g$.
\end{prop}

\begin{prop} Jika $f \in \fml C^\infty(\R^n)$, maka $\left(L_f\right)^{(\bs\alpha)} = L_{f^{(\bs\alpha)}}$ untuk setiap multi-indeks $\bs\alpha$.
\end{prop}

\begin{exer} Jika $\delta$ bukan suatu fungsi pada $\R$, lalu apa yang dimaksud orang ketika mereka menulis
   \[ \lim_{k \sto \infty} \frac{k}{\pi(1+k^2x^2)} = \delta(x)\quad? \]
Tunjukkan bahwa, jika ditafsirkan dengan tepat (sebagai suatu pernyataan tentang distribusi), ungkapan itu bahkan benar. \emph{Petunjuk.}
Jika $s_k(x) = k\,\pi^{-1}(1+k^2x^2)^{-1}$ dan $\phi$ suatu fungsi uji pada $\R$, maka
$\int_{-\infty}^\infty s_k\phi = \int_{-\infty}^\infty s_k\phi(0) + \int_{-\infty}^\infty s_k\eta$ dengan $\eta(x) = \phi(x) - \phi(0)$.
Tunjukkan bahwa $\int_{-\infty}^\infty s_k\eta \sto 0$ ketika $k \sto \infty$.
\end{exer}

\begin{exer} Misalkan $f_n(x) = \sin nx$ untuk $n \in \N$ dan $x \in \R$. Barisan $(f_n)$ tidak konvergen secara klasik;
tetapi barisan itu konvergen secara distribusional (ke $0$). Nyatakan pernyataan ini dengan tepat, lalu buktikan. \emph{Petunjuk.} Carilah antiturunannya.
\end{exer}

\begin{prop} Jika $\phi$ suatu fungsi mulus pada $\R$ dengan tumpuan kompak, maka
    \[ \lim_{n \sto \infty}\frac2\pi \int_{-\infty}^\infty \frac{n^3x}{\bigl(1 + n^2x^2\bigr)^2}\,\, \phi(x)\,dx = \phi'(0)\,.\]
\end{prop}

\begin{exer} Berikan makna pada ungkapan
   \[ 2\pi\sum_{n = -\infty}^\infty \delta(x-2\pi n) =  1 +  2\sum_{n = 1}^\infty \cos nx\,, \]
dan tunjukkan bahwa, jika ditafsirkan dengan tepat, ungkapan itu benar. \emph{Petunjuk.} Siapa pun yang cukup bejat untuk percaya bahwa $\delta(x)$ \emph{berarti}
sesuatu kemungkinan akan menafsirkan $\delta(x-a)$ sebagai hal yang sama dengan $\delta_a(x)$. Anda dapat memulai proses
merasionalisasikannya dengan menghitung deret Fourier fungsi $g$ yang didefinisikan oleh $g(x) = \frac1{4\pi}x^2 - \frac12 x$ pada $[0,2\pi]$
dan kemudian diperluas secara periodik ke $\R$. Anda mungkin perlu mencari teorema tentang konvergensi titik demi titik dan seragam dari deret Fourier
serta tentang pendiferensialan suku demi suku deret semacam itu.
\end{exer}

\begin{defn} Misalkan $\Omega$ suatu subhimpunan terbuka dari $\R^n$, $u \in \fml D^*(\Omega)$, dan $f \in \fml C^\infty(\Omega)$.
 \index{perkalian!distribusi dengan fungsi mulus}%
 \index{<binopconcat@$fu$ (hasil kali fungsi mulus dengan distribusi)}%
Definisikan
   \[ (fu)(\phi) := u(f\phi) \]
untuk semua fungsi uji $\phi$ pada~$\Omega$. Secara ekuivalen, kita dapat menulis $\langle fu,\phi \rangle = \langle u,f\phi \rangle$.
\end{defn}

\begin{prop} Fungsional $fu$ dalam definisi sebelumnya merupakan distribusi pada~$\Omega$.
\end{prop}

\begin{proof} Lihat~\cite{Rudin:1991}, halaman 159, Seksi 6.15.    \ns
\end{proof}






















\section{Konvolusi}

Seksi ini dan Seksi~\ref{section_Fourier_transform} memberikan pengantar yang \emph{sangat} singkat tentang konvolusi
dan transformasi Fourier distribusi. Untuk banyak pembuktian, serta uraian yang lebih terperinci, pembaca dirujuk
ke buku \emph{Functional Analysis} karya Walter Rudin yang elegan~\cite{Rudin:1991}.

Untuk menghindari munculnya terlalu banyak faktor berbentuk $\sqrt{2\pi}$ beserta kebalikannya dalam rumus-rumus, selama
sisa bab ini kita akan mengintegralkan terhadap
 \index{ukuran Lebesgue ternormalisasi}%
 \index{m@$m$ (ukuran Lebesgue ternormalisasi)}%
\df{ukuran Lebesgue ternormalisasi} $m$ pada~$\R$. Ukuran ini didefinisikan oleh
   \[ m(A) := \frac1{\sqrt{2\pi}}\lambda(A) \]
untuk setiap subhimpunan terukur Lebesgue $A$ dari $\R$ (dengan $\lambda$ ukuran Lebesgue biasa pada~$\R$).

Kita meninjau kembali beberapa fakta dasar dari analisis real mengenai konvolusi fungsi bernilai skalar. Ingat bahwa suatu fungsi bernilai kompleks atau
real diperluas pada $\R$ disebut
 \index{terintegralkan Lebesgue}%
 \index{terintegralkan}%
\emph{terintegralkan Lebesgue} jika $\int_\R\abs f\, dm < \infty$
(dengan $m$ ukuran Lebesgue ternormalisasi pada~$\R$).
 \index{L@$\lfs 1(\R)$!ruang fungsi terintegralkan Lebesgue}%
 \index{L@$\lfs 1(\R)$!sebagai ruang Banach}%
 \index{Banach!ruang!$\lfs 1(\R)$ sebagai}%
Dengan $\lfs 1(\R)$ kita nyatakan ruang Banach semua kelas ekuivalensi fungsi terintegralkan Lebesgue pada~$\R$; dua fungsi \emph{ekuivalen}
jika ukuran Lebesgue himpunan tempat keduanya berbeda adalah nol. Norma pada ruang Banach ini
   \index{<unopn@$\norm{\hphantom{x}}_1$ (norma-$1$)}%
diberikan oleh
   \[ {\norm f}_1 := \int_\R \abs f\,dm. \]

\begin{prop} Misalkan $f$ dan $g$ fungsi-fungsi terintegralkan Lebesgue pada~$\R$. Maka fungsi
   \[ y \mapsto f(x - y)g(y) \]
termasuk dalam $\lfs 1(\R)$ untuk hampir setiap~$x$. Definisikan suatu
 \index{<binopast@$f \ast g$ (konvolusi dua fungsi)}%
fungsi $f \ast g$ dengan
   \[ (f \ast g)(x) := \int_\R f(x - y)g(y)\,dm(y) \]
pada setiap titik tempat ruas kanan terdefinisi. Maka fungsi $f \ast g$ termasuk dalam $\lfs 1(\R)$.
\end{prop}

\begin{proof} Lihat~\cite{HewittS:1965}, Teorema 21.31.  \ns
\end{proof}

\begin{defn} Fungsi $f \ast g$ yang didefinisikan di atas adalah
 \index{konvolusi!fungsi-fungsi dalam $\lfs 1(\R)$}%
\df{konvolusi} dari $f$ dan~$g$.
\end{defn}

\begin{prop} Ruang Banach $\lfs 1(\R)$ dengan konvolusi merupakan aljabar Banach komutatif. Namun, aljabar ini \emph{tidak} beridentitas.
\end{prop}

\begin{proof} Lihat~\cite{HewittS:1965}, Teorema 21.34 dan 21.35.  \ns
\end{proof}

\begin{defn}\label{defn_Fourier_transform} Untuk setiap $f$ dalam $\lfs 1(\R)$, definisikan suatu fungsi $\hat f$ dengan
    \[ \hat f(x) = \int_{-\infty}^\infty f(t) e^{-itx}\,dm(t). \]
Fungsi $\hat f$ adalah
 \index{Fourier!transformasi}%
 \index{transformasi!Fourier}%
 \index{F@$\mathfrak F(f)$ (transformasi Fourier dari~$f$)}%
 \index{<unoptop@$\hat f$ (transformasi Fourier dari~$f$)}%
\df{transformasi Fourier} dari~$f$. Kadang-kadang kita menulis $\mathfrak Ff$ untuk~$\hat f$.
\end{defn}

\begin{thm}[Lemma Riemann--Lebesgue] Jika $f \in \lfs 1(\R)$,
 \index{lemma Riemann--Lebesgue}%
maka $\hat f \in \fml C_0(\R)$.
\end{thm}

\begin{proof} Lihat~\cite{HewittS:1965}, 21.39.   \ns
\end{proof}

\begin{prop} Transformasi Fourier
   \[ \mathfrak F\colon \lfs 1(\R) \sto \fml C_0(\R)\colon f \mapsto \hat f \]
merupakan homomorfisme aljabar Banach. Kedua aljabar itu tidak beridentitas.
\end{prop}

Untuk penerapan, segi paling berguna dari proposisi sebelumnya ialah bahwa transformasi Fourier
mengubah konvolusi dalam $\lfs 1(\R)$ menjadi
perkalian titik demi titik dalam~$\fml C_0(\R)$; yakni, $\wh{f\ast g} =\hat f\hat g$.

Berikutnya kita mendefinisikan apa yang dimaksud dengan mengambil konvolusi suatu distribusi dan suatu fungsi uji.

\begin{notn} Untuk suatu titik $x \in \R$ dan suatu fungsi bernilai skalar $\phi$ pada $\R$
 \index{phi@$\phi_x$}%
definisikan
   \[ \phi_x\colon \R \sto \K\colon y \mapsto \phi(x - y). \]
\end{notn}

\begin{prop} Jika $\phi$ dan $\psi$ merupakan fungsi-fungsi bernilai skalar pada $\R$ dan $x \in \R$, maka
   \[ (\phi + \psi)_x = \phi_x + \psi_x\,. \]
\end{prop}

\begin{prop}\label{prop_convol_distr1} Untuk $u \in \fml D^*(\R)$ dan $\phi \in \fml D(\R)$,
 \index{<binopast@$u \ast \phi$ (konvolusi distribusi dan fungsi uji)}%
tetapkan
   \[ (u \ast \phi)(x) := \langle u,\phi_x \rangle \]
untuk semua $x \in \R$. Maka $u \ast \phi \in \fml C^\infty(\R)$.
\end{prop}

\begin{proof} Lihat~\cite{Rudin:1991}, Teorema 6.30(b).  \ns
\end{proof}

\begin{defn} Ketika $u$ dan $\phi$ seperti dalam proposisi sebelumnya, fungsi $u \ast \phi$ adalah
 \index{konvolusi!distribusi dan fungsi uji}%
\df{konvolusi} dari $u$ dan~$\phi$.
\end{defn}

\begin{prop} Jika $u$ dan $v$ merupakan distribusi pada $\R$ dan $\phi$ suatu fungsi uji pada $\R$, maka
   \[ (u + v) \ast \phi = (u \ast \phi) + (v \ast \phi)\,. \]
\end{prop}

\begin{prop} Jika $u$ suatu distribusi pada $\R$ dan $\phi$ serta $\psi$ fungsi-fungsi uji pada $\R$, maka
   \[ u \ast (\phi + \psi) = (u \ast \phi) + (u \ast \psi)\,. \]
\end{prop}

\begin{exam} Jika $H$ fungsi Heaviside dan $\phi$ suatu fungsi uji pada~$\R$, maka
   \[ (\wt H \ast \phi)(x) = \int_{-\infty}^x \phi(t)\,dt. \]
\end{exam}

\begin{prop}\label{prop_convol_distr2} Jika $u \in \fml D^*(\R)$ dan $\phi \in \fml D(\R)$, maka
   \[ D^{\bs\alpha}(u \ast \phi) = (D^{\bs\alpha}u) \ast \phi = u \ast \bigl(D^{\bs\alpha}(\phi)\bigr). \]
\end{prop}

\begin{proof} Lihat~\cite{Rudin:1991}, Teorema 6.30(b).  \ns
\end{proof}

\begin{prop}\label{prop_convol_distr3} Jika $u \in \fml D^*(\R)$ dan $\phi$, $\psi \in \fml D(\R)$, maka
   \[ u \ast (\phi \ast \psi) = (u \ast \phi) \ast \psi. \]
\end{prop}

\begin{proof} Lihat~\cite{Rudin:1991}, Teorema 6.30(c).  \ns
\end{proof}

\begin{prop} Misalkan $u$, $v \in \fml D^*(\R)$. Jika
   \[ u \ast \phi = v \ast \phi \]
untuk semua $\phi \in \fml D(\R)$, maka $ u = v$.
\end{prop}

\begin{defn} Misalkan $\Omega$ suatu subhimpunan terbuka dari $\R^n$ dan $u \in \fml D^*(\R^n)$. Kita mengatakan bahwa ``$u = 0$ pada $\Omega$'' jika $u(\phi) = 0$
untuk setiap $\phi \in \fml D(\Omega)$. Dengan semangat yang sama, kita mengatakan bahwa ``dua distribusi $u$ dan $v$ sama pada $\Omega$'' jika $u - v = 0$
pada~$\Omega$. Suatu titik $x \in \R^n$ termasuk dalam
 \index{tumpuan!distribusi}%
 \index{supp@$\supp u$ (tumpuan distribusi~$u$)}%
\df{tumpuan} dari $u$ jika \emph{tidak ada} lingkungan terbuka bagi $x$ yang padanya $u = 0$. Nyatakan tumpuan $u$ dengan~$\supp u$.
\end{defn}

\begin{exam} Tumpuan distribusi delta Dirac $\delta$ adalah~$\{0\}$.
\end{exam}

\begin{exam} Tumpuan distribusi Heaviside $\wt H$ adalah~$[0,\infty)$.
\end{exam}

\begin{rem} Andaikan suatu distribusi $u$ pada $\R^n$ memiliki tumpuan kompak. Maka dapat ditunjukkan bahwa $u$ memiliki perluasan unik menjadi
fungsional linear kontinu pada~$\fml C^\infty(\R^n)$, yakni keluarga semua fungsi mulus pada~$\R^n$. Dalam kondisi ini, kesimpulan
Proposisi~\ref{prop_convol_distr1} tetap benar dan definisi berikut berlaku. Jika $u \in \fml D^*(\R^n)$ memiliki tumpuan kompak
dan $\phi \in \fml C^\infty(\R^n)$, maka kesimpulan Proposisi~\ref{prop_convol_distr2} tetap benar. Selain itu, jika $\psi$
merupakan fungsi uji pada $\R^n$, maka $u \ast \psi$ merupakan fungsi uji pada $\R^n$ dan kesimpulan~\ref{prop_convol_distr3} tetap berlaku.
Untuk pembahasan dan verifikasi menyeluruh atas pernyataan-pernyataan ini, lihat~\cite{Rudin:1991}, Teorema 6.24 dan~6.35.
\end{rem}

\begin{prop} Jika $u$ dan $v$ merupakan distribusi pada $\R^n$ dan paling sedikit salah satunya memiliki tumpuan kompak, maka terdapat distribusi tunggal
$u \ast v$ sedemikian sehingga
   \[ (u \ast v) \ast \phi = u \ast (v \ast \phi) \]
untuk semua fungsi uji $\phi$ pada~$\R^n$.
\end{prop}

\begin{proof} Lihat \cite{Rudin:1991}, Definisi 6.36\emph{ dan seterusnya}.   \ns
\end{proof}

\begin{defn} Dalam proposisi sebelumnya, distribusi $u \ast v$ adalah
 \index{konvolusi!dua distribusi}%
 \index{<binopast@$u \ast v$ (konvolusi dua distribusi)}%
\df{konvolusi} dari $u$ dan~$v$.
\end{defn}

\begin{exam} Jika $\delta$ distribusi delta Dirac, maka
   \[ \tilde 1 \ast \delta\,' = \tilde 0\,. \]
\end{exam}

\begin{exam} Jika $\delta$ distribusi delta Dirac dan $H$ fungsi Heaviside, maka
   \[ \delta\,' \ast \wt H = \delta\,. \]
\end{exam}

\begin{prop} Jika $u$ dan $v$ merupakan distribusi pada $\R^n$ dan paling sedikit salah satunya memiliki tumpuan kompak, maka $u \ast v = v \ast u$.
\end{prop}

\begin{proof} Lihat \cite{Rudin:1991}, Teorema 6.37(a).  \ns
\end{proof}

\begin{exer} Buktikan proposisi sebelumnya dengan asumsi bahwa $u$ dan $v$ \emph{keduanya} memiliki tumpuan kompak. \emph{Petunjuk.}
Gunakan $(u \ast v) \ast (\phi \ast \psi)$, dengan $\phi$ dan $\psi$ fungsi-fungsi uji, sambil mengingat bahwa konvolusi fungsi
bersifat komutatif.
\end{exer}

\begin{prop} Jika $u$ dan $v$ merupakan distribusi pada $\R^n$ dan paling sedikit salah satunya memiliki tumpuan kompak, maka
  \[ \supp(u \ast v) \subseteq \supp u + \supp v. \]
\end{prop}

\begin{proof} Lihat \cite{Rudin:1991}, Teorema 6.37(b).  \ns
\end{proof}

\begin{prop}\label{prop_assoc_convol_distr} Jika $u$, $v$, dan $w$ merupakan distribusi pada $\R^n$ dan paling sedikit dua di antaranya memiliki
tumpuan kompak, maka
    \[ u \ast (v \ast w) = (u \ast v) \ast w. \]
\end{prop}

\begin{exam} Ungkapan $\tilde 1 \ast \delta\,' \ast \wt H$ bersifat ambigu. Jadi, hipotesis mengenai tumpuan
distribusi-distribusi dalam proposisi sebelumnya tidak dapat dihilangkan.
\end{exam}

\begin{prop} Terhadap penjumlahan, perkalian skalar, dan konvolusi, keluarga semua distribusi pada $\R$ yang bertumpuan kompak
merupakan aljabar komutatif beridentitas.
\end{prop}













\section{Solusi Distribusional untuk Persamaan Diferensial Biasa}

\begin{notn} Misalkan $\Omega$ suatu subhimpunan terbuka tak kosong dari $\R$, $n \in \N$, dan $a_0$, $a_1$, \dots $a_n \in \fml C^\infty(\Omega)$.
Kita meninjau operator diferensial (biasa)
 \index{diferensial!operator}%
 \index{operator!diferensial}%
berikut:
   \[ L = \sum_{k=0}^n a_k(x)\frac{d^k}{dx^k} \]
beserta persamaan diferensial (biasa)
 \index{diferensial!persamaan}%
 \index{persamaan!diferensial}%
yang terkait dengannya,
   \begin{equation}\label{eqn_ODE1}
        Lu = v
   \end{equation}
dengan $u$ dan $v$ distribusi-distribusi pada~$\Omega$. Sehubungan dengan Persamaan~\eqref{eqn_ODE1}, kita juga dapat meninjau dua variasi:
persamaan
   \begin{equation}\label{eqn_ODE2}
        \langle Lu, \phi\rangle = \langle v, \phi\rangle
   \end{equation}
dengan $\phi$ suatu fungsi uji, dan persamaan
   \begin{equation}\label{eqn_ODE3}
        \langle u, L^*\phi\rangle = \langle v, \phi\rangle
   \end{equation}
dengan $\phi$ suatu fungsi uji dan $L^* := \displaystyle\sum_{k=0}^n (-1)^k \frac{d^k(a_k\phi)}{dx^k}$ merupakan
 \index{adjoin!formal}%
 \index{adjoin formal}%
\df{adjoin formal} dari~$L$.
\end{notn}

\begin{defn} Andaikan bahwa dalam persamaan-persamaan di atas distribusi $v$ diberikan. Jika terdapat distribusi regular $u$ yang
memenuhi~\eqref{eqn_ODE1}, kita mengatakan bahwa $u$ merupakan solusi
 \index{solusi klasik}%
 \index{solusi!klasik}%
\df{klasik} bagi persamaan diferensial itu. Jika $u$ merupakan distribusi regular yang bersesuaian dengan suatu fungsi yang
\emph{tidak} memenuhi~\eqref{eqn_ODE1} dalam pengertian klasik, tetapi \emph{memenuhi}~\eqref{eqn_ODE2} untuk setiap fungsi uji~$\phi$,
maka kita mengatakan bahwa $u$ merupakan solusi
 \index{lemah!solusi}%
 \index{solusi!lemah}%
\df{lemah} bagi persamaan diferensial itu. Selanjutnya, jika $u$ merupakan distribusi singular yang memenuhi~\eqref{eqn_ODE3} untuk setiap fungsi
uji~$\phi$, kita mengatakan bahwa $u$ merupakan solusi
 \index{solusi distribusional}%
 \index{solusi!distribusional}%
\df{distribusional} bagi persamaan diferensial itu. Keluarga solusi
 \index{solusi diperumum}%
 \index{solusi!diperumum}%
\df{diperumum} mencakup semua solusi klasik, lemah, dan distribusional.
\end{defn}

\begin{exam} Solusi diperumum bagi persamaan $\dfrac{du}{dx} = 0$ pada~$\R$ hanyalah distribusi konstan (yakni,
distribusi regular yang berasal dari fungsi konstan). Jadi, setiap solusi diperumum merupakan solusi klasik.
\end{exam}

\begin{proof}[\emph{Petunjuk untuk bukti}] Andaikan $u$ suatu solusi diperumum bagi persamaan itu. Pilih suatu fungsi uji $\sbsb\phi 0$ sedemikian
sehingga $\int_\R \sbsb\phi 0 = 1$. Misalkan $\phi$ suatu fungsi uji sebarang pada~$\R$. Definisikan
   \[ \psi(x) = \phi(x) - \sbsb\phi 0(x)\int_\R \phi \]
untuk semua $x \in \R$. Amati bahwa $\langle u,\psi \rangle = 0$, lalu hitung $\langle u,\phi \rangle$.  \ns
\end{proof}

\begin{exam} Distribusi Heaviside adalah solusi lemah persamaan diferensial $x\dfrac{du}{dx}~=~0$ pada~$\R$. Tentu saja,
distribusi-distribusi konstan merupakan solusi klasik.
\end{exam}

\begin{exam} Distribusi Dirac $\delta$ merupakan solusi distribusional bagi persamaan diferensial $x^2\dfrac{du}{dx} = 0$ pada~$\R$.
(Distribusi Heaviside merupakan solusi lemah; dan distribusi-distribusi konstan merupakan solusi klasik.)
\end{exam}

\begin{exam} Untuk $p = 0,1,2,\dots$, turunan ke-$p$, $\dfrac{d^p\delta}{dx^p}$, dari distribusi delta Dirac merupakan solusi
bagi persamaan diferensial $x^{p+2}\dfrac{du}{dx} = 0$ pada~$\R$.
\end{exam}

\begin{exer} Carilah suatu solusi distribusional bagi persamaan diferensial $x \dfrac{du}{dx} + u = 0$ pada~$\R$.
\end{exer}










\section{Transformasi Fourier}\label{section_Fourier_transform}

Dalam~\ref{defn_Fourier_transform}, kita mendefinisikan transformasi Fourier suatu fungsi $\lfs 1\R$. Cara ekuivalen untuk menyatakan
definisi ini adalah melalui konvolusi:
  \[ (\mathfrak F f)(x) = (f \ast \varepsilon^x)(0) \]
untuk $x \in \R$, dengan $\varepsilon^x$ fungsi $t \mapsto e^{itx}$.

\begin{defn} Kita ingin memperluas transformasi yang berguna ini ke distribusi. Namun, ternyata ruang $\fml D^*(\R)$ terlampau besar
untuk mendukung definisi yang wajar baginya. Dengan memulai dari himpunan ``fungsi uji'' yang lebih besar, yakni fungsi-fungsi dalam ruang Schwartz
$\fml S(\R)$ (yang memuat $\fml D(\R)$ sebagai subruang rapat), kita memperoleh ruang fungsional linear kontinu yang lebih kecil; pada ruang ini, untungnya, kita
dapat mendefinisikan suatu \emph{transformasi Fourier} secara alami.
\end{defn}

Sebelum mendefinisikan transformasi Fourier pada distribusi, mari kita ingat kembali dua fakta penting dari analisis real mengenai
transformasi Fourier yang bekerja pada fungsi mulus terintegralkan.

\begin{prop}[Rumus Inversi Fourier] Jika $f \in \fml S(\R)$,
 \index{Fourier!rumus inversi}%
maka
   \[ f(x) = \int_\R \hat f \varepsilon^x\,dm\,. \]
\end{prop}

\begin{prop}\label{prop_FT_isomor1} Transformasi Fourier
   \[ \mathfrak F\colon \fml S(\R) \sto \fml S(\R)\colon f \mapsto \hat f \]
merupakan isomorfisme ruang Fr\'echet (yakni, pemetaan linear kontinu bijektif dengan invers kontinu), dan untuk semua
$f \in \fml S(\R)$ serta $x \in \R$, kita mempunyai $\mathfrak F^2f(x) = f(-x)$, sehingga $\mathfrak F^4f = f$.
\end{prop}

\begin{proof} Pembuktian kedua hasil sebelumnya, bersama banyak informasi lain tentang transformasi yang menarik ini
dan beragam penerapannya pada persamaan diferensial biasa maupun parsial, dapat ditemukan di berbagai sumber.
Beberapa yang segera terlintas adalah~\cite{Cerda:2010}, Bab~7; \cite{Horvath:1966}, Bab~4, Seksi~11;
\cite{McDonaldW:1999}, Seksi~11.3; \cite{Rudin:1991}, Bab~7; \cite{Treves:1967}, Bab~25; dan \cite{Yosida:1965},
Bab~VI.   \ns
\end{proof}

\begin{defn} Fungsi inklusi $\iota\colon \fml D(\R) \sto \fml S(\R)$ bersifat kontinu. Maka pemetaan antara ruang-ruang dualnya
   \[ u\colon \fml S^*(\R) \sto \fml D^*(\R)\colon \tau \mapsto u_\tau = \tau \circ \iota \]
merupakan isomorfisme ruang vektor dari $\fml S^*(\R)$ ke suatu ruang distribusi pada~$\R$, yang kita sebut ruang
 \index{distribusi tempered}%
\df{distribusi tempered} (sebagian penulis memilih frasa
 \index{distribusi!temperate}%
\df{distribusi temperate}). Lazimnya, fungsional linear $\tau$ dan $u_\tau$ di atas diidentifikasi. Jadi, $\fml S^*(\R)$
sendiri dipandang sebagai ruang distribusi tempered.
\end{defn}

Jika Anda kesulitan memverifikasi salah satu pernyataan dalam definisi sebelumnya, lihat Teorema 7.10\emph{ dan seterusnya} dalam~\cite{Rudin:1991}.

\begin{exam} Setiap distribusi dengan tumpuan kompak merupakan distribusi tempered.
\end{exam}

\begin{defn} Definisikan
 \index{Fourier!transformasi}%
 \index{transformasi!Fourier}%
 \index{F@$\mathfrak F(u)$ (transformasi Fourier dari~$u$)}%
 \index{<unoptop@$\hat u$ (transformasi Fourier dari~$u$)}%
\df{transformasi Fourier} $\hat u$ dari suatu distribusi tempered $u \in \fml S^*(\R)$ dengan
   \[ \hat u(\phi) = u(\hat\phi) \]
untuk semua $\phi \in \fml S(\R)$. Notasi $\mathfrak Fu$ merupakan alternatif bagi~$\hat u$.
\end{defn}

\begin{prop} Jika $f \in \lfs1\R$, maka kita dapat memperlakukan $\tilde f$ sebagai distribusi tempered dan, dalam hal ini,
   \[ \hat{\tilde f} = \tilde{\hat f}\,. \]
\end{prop}

\begin{exam} Transformasi Fourier fungsi delta Dirac diberikan oleh
   \[ \hat\delta = \tilde{\vc 1}\,. \]
\end{exam}

\begin{exam} Transformasi Fourier distribusi regular $\tilde{\vc 1}$ diberikan oleh
   \[ \hat{\tilde{\vc 1}} = \delta\,. \]
\end{exam}

Proposisi~\ref{prop_FT_isomor1} memiliki generalisasi yang memuaskan bagi distribusi tempered.

\begin{prop} Transformasi Fourier $\mathfrak F\colon \fml S^*(\R) \sto \fml S^*(\R)$ merupakan pemetaan linear kontinu
bijektif yang inversnya juga kontinu. Lebih lanjut, $\mathfrak F^4$ merupakan pemetaan identitas pada~$\fml S^*(\R)$.
\end{prop}

Kita akan kembali sejenak ke pembahasan transformasi Fourier dalam bab berikutnya.
















\endinput
